\textbf{Vision-Language Models (VLMs).} We consider VLMs that pair a Vision Transformer (ViT)~\citep{dosovitskiy2020vit} with an autoregressive text decoder. The ViT encodes an input image $\mathbf{I}$ into a sequence of patch embeddings $\{v_i\}_{i=1}^{N}$, which serve as \emph{patch tokens}. A projection module maps these tokens into the language model's embedding space, where they act as conditioning context. The decoder then predicts tokens $\{w_t\}_{t=1}^{T}$ autoregressively under the standard next-token prediction objective. In LLaVA-1.5-7B, the ViT is a CLIP ViT-L/14 operating at $336{\times}336$ resolution, giving an $N = 24{\times}24 = 576$ patch grid, and the decoder is Vicuna-7B.

\textbf{Saliency Maps.} Because patch tokens and text tokens are processed jointly within the decoder's self-attention layers, each generated token $w_t$ attends to all $N$ patch tokens. We form a saliency map $\mathbf{S}_t$ for $w_t$ by aggregating these token-to-patch attention weights across layers and heads and reshaping the result onto the ViT patch grid, producing a spatial map that highlights the regions most relevant to predicting $w_t$. Throughout, we normalize $\mathbf{S}_t$ to a probability distribution over spatial locations so that maps are comparable across images and tokens. This saliency signal is the quantity our objective supervises (\Cref{sec:method}) and our metrics evaluate (\Cref{sec:experiment}).
